
\documentclass[12pt]{article}
\usepackage[margin=1in]{geometry}
\usepackage[utf8]{inputenc}
\usepackage[T1]{fontenc}
\usepackage{hyperref}
\usepackage{enumitem}
\usepackage{xcolor}
\usepackage{titlesec}
\usepackage{fancyhdr}
\usepackage{tcolorbox}
\tcbuselibrary{breakable}
\usepackage{setspace}
\usepackage{parskip}
\setlength{\parskip}{6pt plus 2pt minus 1pt}

% UofSC Color Definitions
\definecolor{garnet}{RGB}{115,0,10}
\definecolor{black}{RGB}{0,0,0}
\definecolor{rose}{RGB}{204,46,64}
\definecolor{atlantic}{RGB}{70,106,159}
\definecolor{congaree}{RGB}{31,65,77}
\definecolor{neutral90}{RGB}{54,54,54}

\hypersetup{
  colorlinks=true,
  urlcolor=atlantic,
  linkcolor=black,
  citecolor=black
}

% Custom compact title
\newcommand{\doctitle}{%
    \noindent{\LARGE\bfseries Daily Gamecock \textit{Opinion}}\\[2pt]
    {\LARGE\bfseries Requesting Statements}\\[4pt]
    {\normalsize J.C. Vaught, Opinion Editor. Fall 2025}\\[-2pt]
    \noindent\rule{\textwidth}{1pt}
    \vspace{8pt}
}

\pagestyle{fancy}
\fancyhf{}
\lhead{Daily Gamecock \textit{Opinion}}
\rhead{Requesting Statements}
\cfoot{\thepage}
\renewcommand{\headrulewidth}{0.4pt}

% Section formatting
\titleformat{\section}
{\normalfont\Large\bfseries\color{garnet}}{\thesection}{1em}{}

\titleformat{\subsection}
{\normalfont\large\bfseries\color{neutral90}}{\thesubsection}{1em}{}

\setlist[itemize]{topsep=3pt,itemsep=3pt,parsep=0pt}
\setlist[enumerate]{topsep=3pt,itemsep=3pt,parsep=0pt}

\newcommand{\key}[1]{\textbf{#1}}
\newcommand{\deadline}[1]{\textcolor{rose}{\textbf{#1}}}
\newcommand{\good}[1]{\textbf{\textcolor{congaree}{GOOD: #1}}}
\newcommand{\bad}[1]{\textbf{\textcolor{rose}{BAD: #1}}}

\begin{document}
\thispagestyle{plain}

\doctitle

\begin{tcolorbox}[colback=neutral90!10,colframe=neutral90,boxrule=1pt,sharp corners,left=8pt,right=8pt,top=8pt,bottom=8pt,title={Purpose \& Scope},fonttitle=\bfseries,coltitle=white,colbacktitle=neutral90,toptitle=4pt,bottomtitle=4pt]
This guide explains \emph{when} to ask for an official statement, \emph{how} to ask by email, and provides templates and examples for contacting universities, government agencies, companies, and other institutions. It is written for the Daily Gamecock \textit{\textbf{Opinion}} section, but the practices apply broadly across student and professional journalism.
\end{tcolorbox}


\section{When you must seek a statement (\emph{right of reply})}
Seeking comment is both an ethical obligation and a practical legal safeguard:
\begin{itemize}
  \item \key{Criticism or allegations}: If an opinion piece criticizes a person or institution, you must ask them to respond before publication. 
  \item \key{Potentially harmful assertions}: If your draft includes claims that could affect someone's reputation (e.g., negligence, incompetence, discrimination), offer a fair chance to respond; this can help demonstrate lack of malice and reduce libel risk.
  \item \key{Policy or decision coverage}: When discussing university or government \emph{actions}, seek on-the-record reasoning, timelines, or corrective steps from the relevant office (PIO/press office, dean, agency spokesperson).
\end{itemize}

\paragraph{When it's optional but smart.} If an opinion piece is primarily praise or personal reflection, comment is optional; however, if you cite new facts or statistics about an institution, confirm those details or offer a chance to clarify.

\section{Who to contact}
\begin{enumerate}
  \item \key{Universities}: Start with the \emph{media relations/communications} office. Your editor should have a contact there.
  \item \key{Government}: Identify the agency's press office or public information officer (PIO); for elected officials, the press secretary. If facts hinge on records, consider a public-records request in parallel (FOIA/state open-records).
  \item \key{Companies/nonprofits}: Use the media relations email or the executive responsible for the policy at issue; many list a press inbox specifically for \emph{requests for comment}.
\end{enumerate}

\section{How to ask (email anatomy)}
\subsection*{Subject line}
Include outlet, topic, and a clear deadline. Concise subjects that signal the outlet and topic help sources triage quickly. Examples: 
\begin{itemize}
  \item \texttt{Request for comment -- USC housing waitlist concerns}
  \item \texttt{Media request (Daily Gamecock) -- Student fees oversight}
\end{itemize}


\subsection*{Opening and context}
\begin{itemize}
  \item Identify yourself and the \emph{Daily Gamecock}.
  \item State the \emph{opinion thesis} you are testing.
  \item List \emph{specific} questions tied to that thesis.
\end{itemize}

\subsection*{Ground rules on attribution}
Default is \key{on the record}. The Daily Gamecock does not accept ``Off the record'' or ``on background'' sources and statements without explicit approval from the managing editors. If a source requests background for safety or sensitivity, confirm the terms with your editors.

\subsection*{Reasonable time \& follow-up}
Give institutions a fair window to respond relative to your publication schedule; document attempts to reach them and note ``declined to comment'' only after bona fide efforts. If time is short, say so explicitly and offer phone availability for faster turnaround.

\section{Email templates (copy, customize, send)}
\subsection*{A. University}
\begin{tcolorbox}[colback=neutral90!5,colframe=neutral90,boxrule=1pt,sharp corners,left=8pt,right=8pt,top=8pt,bottom=8pt]
\textbf{Subject:} Request for comment --- [Topic]\\[4pt]
Hello,\\[2pt]
I'm [Name], an opinion writer with the \emph{Daily Gamecock}. I'm writing about [\emph{thesis in one sentence}].\\[6pt]
To ensure fairness and accuracy, could I get a statement on the following?\\[-4pt]
\begin{enumerate}[label=\arabic*.]
  \item What specific policies or constraints led to [issue] this semester?
  \item What steps, if any, will USC take before [date] to address [concrete fix]?
  \item How should affected students seek help in the meantime (office, link, timeline)?
\end{enumerate}
Thank you,\\
Name\\
\emph{Daily Gamecock} \textbar{} [email] \textbar{} [phone]
\end{tcolorbox}

\subsection*{B. Government agency}
\begin{tcolorbox}[colback=neutral90!5,colframe=neutral90,boxrule=1pt,sharp corners,left=8pt,right=8pt,top=8pt,bottom=8pt]
\textbf{Subject:} Media request (Daily \emph{Gamecock}) --- [Agency]  comment on [policy]\\[4pt]
Hello [Press Officer],\\[2pt]
I'm covering [policy/action] and its impact on [group]. For accuracy, please provide on-the-record responses to:\\[-4pt]
\begin{enumerate}[label=\arabic*.]
  \item What is the agency's stated rationale and legal authority for [action]?
  \item What data informed the decision? Please link to any public reports.
  \item What recourse is available to affected students/residents?
\end{enumerate}

Thank you,\\
Name\\
\emph{Daily Gamecock} \textbar{} [email] \textbar{} [phone]
\end{tcolorbox}

\subsection*{C. Company/nonprofit}
\begin{tcolorbox}[colback=neutral90!5,colframe=neutral90,boxrule=1pt,sharp corners,left=8pt,right=8pt,top=8pt,bottom=8pt]
\textbf{Subject:} Request for comment --- [Company] policy on [issue]\\[4pt]
Hello,\\[2pt]
In a forthcoming column about [issue], I'm asking [Company] to respond on the record to:\\[-4pt]
\begin{enumerate}[label=\arabic*.]
  \item Current policy on [specific practice] as it applies to students.
  \item Any planned changes and the timeline.
  \item Your response to critiques that [one-sentence critique].
\end{enumerate}
Thank you,\\
Name\\
\emph{Daily Gamecock} \textbar{} [email] \textbar{} [phone]
\end{tcolorbox}

\subsection*{D. Short-notice follow-up / voicemail script}
\begin{tcolorbox}[colback=neutral90!5,colframe=neutral90,boxrule=1pt,sharp corners,left=8pt,right=8pt,top=8pt,bottom=8pt]
\textbf{Subject:} Quick follow-up for comment by \deadline{2 p.m. ET today} --- [topic]\\[4pt]
Hi [Name] --- following up on the note below. My print deadline is \deadline{2 p.m. ET today}. Even a brief statement addressing items \#1--\#2 would be appreciated. I'm at [phone] and can read the questions aloud if helpful. Thank you.\\[6pt]
\emph{Voicemail:} ``Hi, this is [Name] with the \emph{Daily Gamecock}. I'm requesting comment on [topic] by [deadline/timezone]. My questions are [one-line list]. Please call [phone] or reply to [email]. Thank you.''
\end{tcolorbox}

\section{Examples (good vs.\ needs work)}
\subsection*{Good}
\begin{tcolorbox}[colback=congaree!10,colframe=congaree,boxrule=1pt,sharp corners,left=8pt,right=8pt,top=8pt,bottom=8pt]
\textbf{Subject:} Request for comment by 4 p.m. ET --- Student fee transparency\\
Identifies outlet and deadline; states thesis; asks concrete, answerable questions; offers phone; clarifies on-the-record.
\end{tcolorbox}

\subsection*{Needs work}
\begin{tcolorbox}[colback=rose!10,colframe=rose,boxrule=1pt,sharp corners,left=8pt,right=8pt,top=8pt,bottom=8pt]
\textbf{Subject:} Question\\
Vague subject; no deadline; asks for ``any comment'' without context; only background questions; no attribution terms.
\end{tcolorbox}

\section{Documentation \& transparency}
\begin{itemize}
  \item Keep a log of \emph{who}, \emph{when}, and \emph{how} you sought comment (email, phone, in-person); save screenshots.
  \item If no response (generally 48hrs), write: ``The university did not respond to questions by publication time.''
\end{itemize}

\section{Quick checklist}
\begin{enumerate}
  \item \key{Is there criticism or an allegation?} If yes, \emph{seek comment}.%
  \item \key{Have you written the thesis clearly} in your request?
  \item \key{Are the questions specific} (who/what/when/how) and answerable?
  \item \key{Did you contact the right person/office} and provide a phone number?
  \item \key{Did you document attempts} and fairly characterize non-responses?
\end{enumerate}

\end{document}
